% ============================================================================
%  Part IX: Reference Implementations and Benchmarks
% ============================================================================
\section{Reference Implementations and Benchmarks}
\label{sec:implementations}

Eight reference implementations accompany this specification: Python (canonical and performance-tuned variants), Go, Rust, TypeScript, C, C++, Java, and a CPython C-extension binding. Each implementation passes a shared test suite covering size, variant byte handling, uniqueness across 100{,}000 IDs, monotonicity within a generator, encoding round-trips, HLC observation, HLC drift bounds, clock-clamp behavior, and decoder validation. This part documents the design choices in each implementation. Benchmark numbers are deliberately \emph{not} embedded in this document: the live source of record for measured throughput, the measurement environment, and the procedure is \texttt{tessera-id/benchmarks/methodology.md}, with per-language captured outputs alongside it in \texttt{tessera-id/benchmarks/}.

\begin{calloutwarning}[On the optimization headroom in the measured numbers]
The reference implementations prioritize \emph{specification fidelity and code clarity} over peak throughput. They are written so that a porter can compare them line-by-line against the pseudocode in Part \ref{sec:spec} and convince themselves the implementation is correct; they are not written to extract every available cycle. Production deployments that need more than the rates captured in \texttt{tessera-id/benchmarks/} have substantial headroom and several well-understood paths to it. None of these are taken in the reference code:
\begin{itemize}
  \item \textbf{Bulk CSPRNG buffering.} The reference Python ``fast'' variant pre-fills a CSPRNG buffer to amortize syscall cost; the other reference implementations call \code{getrandom(2)} (or its equivalent) once per generation. A buffered version typically gains 1.5--3$\times$ on generation throughput, at the cost of a larger per-thread memory footprint and slightly more complex state management.
  \item \textbf{SIMD-accelerated encoding.} The Base64URL and Crockford encoders in the reference code use scalar byte-at-a-time loops. Production-grade Base64 implementations (e.g., the ones inside major language runtimes) use SSE/AVX2/AVX-512 to process 12--24 bytes per instruction, achieving 5--10$\times$ the throughput of the scalar reference.
  \item \textbf{Thread-local generators.} Implementations using a shared mutex-guarded generator (the Python module-level \code{new()}, TypeScript's module-level generator, Java's \code{SynchronizedStandardGenerator} wrapper) pay contention cost under high concurrency. One generator per thread, or per request handler, eliminates that cost.
  \item \textbf{Hardware RNG offload.} The bulk of generation cost is the kernel-mode CSPRNG syscall. \code{RDSEED}/\code{RDRAND} on x86-64, or a vectorized hardware RNG, can drive generation rates above the syscall floor for workloads where the CSPRNG syscall is genuinely the bottleneck.
\end{itemize}
We chose not to fold these into the reference because they obscure the spec-to-code mapping that the references exist to make legible. Treat the captured benchmark outputs as a reproducible floor on what the format can do, not a ceiling. A production deployment that wants 2--5$\times$ the throughput will not need to redesign anything.
\end{calloutwarning}

\subsection{Implementation summaries}

\subsubsection{Python (\code{python/tessera\_id.py}, \code{python/tessera\_id\_fast.py})}

The canonical Python implementation in \code{tessera\_id.py} prioritizes readability over performance. It is a faithful one-to-one translation of the specification's pseudocode, designed to serve as a reference for porters. Module-level functions operate on a singleton thread-safe generator (mutex-guarded). The performance-tuned variant in \code{tessera\_id\_fast.py} pre-fills a CSPRNG buffer to amortize syscall cost and uses bytearray manipulation to skip Python-object overhead in the hot path.

\subsubsection{Go (\code{go/tessera\_id.go})}

The Go implementation defines an explicit \code{StandardGenerator} struct and recommends one generator per goroutine to avoid mutex contention. \code{HLCGenerator} is internally locked, as required for correctness. Sortedness is enforced at the test level: the test suite generates 10{,}000 consecutive IDs and asserts that \code{bytes.Compare} returns positive for every consecutive pair.

\subsubsection{Rust (\code{rust/src/lib.rs})}

The Rust implementation uses Edition 2024 syntax and the \code{getrandom} 0.3 crate. The \code{TesseraId} type is a fixed-size struct with derived ordering traits, allowing TesseraIDs to be used directly in \code{BTreeMap} keys and sorted vectors with no allocation overhead. Build with \code{--release} for maximum optimization.

\subsubsection{TypeScript (\code{typescript/tessera\_id.ts})}

The TypeScript implementation targets Node 20+ and modern browsers (relying on Web Crypto's \code{getRandomValues}). It pays close attention to JavaScript number-precision pitfalls: the 48-bit timestamp is decomposed into hi/lo 24-bit halves to avoid \code{Number}'s 53-bit precision floor, and the encoding routines use \code{Uint8Array} throughout.

\subsubsection{C (\code{c/tessera\_id.c}, \code{c/tessera\_id.h})}

The C reference is a single .c/.h pair, no dynamic allocation, suitable for embedded use. It uses \code{getrandom(2)} on Linux and is portable to any POSIX system that provides \code{/dev/urandom}. Built with \code{gcc -O3 -std=c11}. Even with the scalar loop, the Base64URL encoder runs tens of nanoseconds per call on current x86-64 --- the encoding itself is essentially free next to the CSPRNG syscall (measured rates: \texttt{benchmarks/}).

\subsubsection{C++ (\code{cpp/tessera\_id.hpp})}

The C++ reference is header-only, C++20, and adds strong typing on top of the C-level wire format. \code{TesseraId} is a struct with \code{operator<=>}, \code{StandardGenerator} and \code{HLCGenerator} are RAII types, and the HLC generator is \code{std::mutex}-protected. Built with \code{g++ -O3 -std=c++20}.

\subsubsection{Java (\code{java/src/main/java/com/tessera/id/})}

The Java reference targets OpenJDK 25 LTS. \code{TesseraId} is an immutable value class with \code{compareTo} delegating to \code{Arrays.compareUnsigned} on the wire bytes, allowing TesseraIDs to drop into \code{TreeMap} or \code{Comparator} chains directly. \code{StandardGenerator} is \emph{unsynchronized by default}, matching the cross-language locking contract (one generator per thread, or the caller wraps it); a \code{SynchronizedStandardGenerator} wrapper is provided for shared-instance use (TID-00019). \code{HLCGenerator} is always serialized (the HLC algorithm requires it) and surfaces an \code{HLCDriftError} when an observation exceeds the 60-second drift bound; \code{observe} accepts a decoded \code{TesseraId} overload in addition to the raw HLC prefix, so observations are validated before they touch generator state. The codec is a hand-rolled Base64URL encoder/decoder (the JDK's \code{java.util.Base64} pads and uses a different alphabet) plus a Crockford Base32 implementation with case-insensitive decoding and confusable canonicalization.

\subsubsection{CPython extension (\code{cpython/tessera\_id\_ext.c})}

The CPython extension wraps the C reference implementation as a Python module, exposing \code{new()}, \code{new\_hlc()}, \code{encode\_b64()}, \code{decode\_b64()}, \code{encode\_b32()}, \code{decode\_b32()}, \code{validate()}, and field accessors. \code{validate()} binds the C reference's \code{tessera\_id\_validate}, and both string decoders enforce the variant byte at decode time per §\ref{sec:decode-variant-contract} (TID-00020). Built via \code{python3 setup.py build\_ext --inplace}. The extension exists to demonstrate the Python-interpreter-overhead contribution relative to the pure-Python variants; see \texttt{benchmarks/} for measured ratios.

\subsection{Cross-implementation benchmark results}

Earlier revisions of this document embedded a point-in-time results table (per-language ops/sec for generation, encoding, and decoding). That table has been removed: embedded numbers go stale the moment any implementation changes, and a static table cannot carry its own provenance. The live source of record is the repository:

\begin{itemize}
  \item \texttt{tessera-id/benchmarks/methodology.md} --- the measurement environment, toolchain pins, warmup/iteration protocol, and the provenance rules every published number must satisfy.
  \item \texttt{tessera-id/benchmarks/} --- one captured output file per language, regenerated (never hand-edited) by each language's own benchmark harness (see the per-language \texttt{bench} targets in \texttt{tessera-id/CLAUDE.md}).
\end{itemize}

Readers who need current figures should consult those files or re-run the harnesses; the qualitative observations below are properties of the design and hold across the environments measured to date.

\subsection{Observations}

\textbf{Generation is CSPRNG-bound.} Across the native implementations (C, C++, Rust, Go, Java), Standard generation lands within a small factor of each other. The differences come down to syscall wrapper thinness; the work outside the syscall is essentially identical and well below the syscall cost. Go's edge over the others is its lighter-weight \code{getrandom} integration in the runtime.

\textbf{HLC is faster than Standard.} Counterintuitively, HLC generation measures consistently \emph{faster} than Standard generation in the native implementations. The reason is that the Standard generator does an extra random-byte derivation for the counter step (\code{1 + (rnd[0] \& 0x0F)}) on every same-millisecond call, which adds a few CPU cycles compared to HLC's simpler increment. The CSPRNG fetch is the same size in both cases.

\textbf{Java tracks the native pack.} HotSpot's JIT brings Java generation close to C/Rust on Standard generation. The encoder is slower than C or Rust but faster than Node's TypeScript --- the JIT-compiled scalar loop holds up well against hand-written native scalar code.

\textbf{Encoding is essentially free in native code.} The C scalar encoders run at tens of nanoseconds per call --- so fast that for any practical use the encoder is never the bottleneck; whatever consumes the encoded string (network I/O, JSON serialization, log writes) dominates by orders of magnitude. As noted in the optimization caveat above, a SIMD-accelerated encoder would push these rates another 5--10$\times$ higher, but the reference uses the scalar version for clarity.

\textbf{Python's pure-Python overhead is real but predictable.} The canonical Python implementation runs several times slower than C; the CPython extension closes most of that gap, demonstrating that the Python-language overhead --- attribute lookups, integer boxing, function call frames --- is the dominant cost when the implementation is in pure Python.

\textbf{TypeScript hits a JavaScript-runtime floor.} The Node.js Web Crypto \code{getRandomValues} call is the floor; TypeScript's syntactic overhead on top is small. For TS workloads where that floor is insufficient, calling out to a native module is the standard remedy.

\textbf{No implementation is the bottleneck for any realistic workload.} A web service that issues TesseraIDs as part of every request can comfortably handle hundreds of thousands of requests per second on a single core; the rest of the request handling (JSON parsing, database I/O, business logic) will dominate by orders of magnitude before TesseraID generation does.
